\subsection{Distributions}

We keep the notations of the preceding number; U therefore denotes an open subset of \(R^{n}\) .

DEFINITION 1. --- The dual space of \(\mathcal{D}(\mathrm{U})\) , endowed with the topology of bounded convergence, is called the space of distributions on U. It is denoted by \(\mathcal{D}'(\mathrm{U})\) .

A distribution on U is therefore a continuous linear form on \(\mathcal{D}(\mathrm{U})\) .

If \(f\) is a linear form on \(\mathcal{D}(\mathrm{U})\) (and in particular if \(f\) is a distribution) and if \(\varphi\) belongs to \(\mathcal{D}(\mathrm{U})\) , we denote by \(\langle f, \varphi \rangle\) the evaluation of \(f\) at \(\varphi\) .

Since \(\mathcal{D}(\mathrm{U})\) is bornological, the space \(\mathcal{D}'(\mathrm{U})\) is complete (TVS, III, p. 24, cor. 1). Since \(\mathcal{D}(\mathrm{U})\) is a Montel space, the same holds for \(\mathcal{D}'(\mathrm{U})\) (EVT, IV, p. 19, prop. 9). In particular, the space \(\mathcal{D}'(\mathrm{U})\) is reflexive (EVT, IV, p. 16, th. 2).

Lemma 4. --- Let f be a linear map from \(\mathcal{D}(\mathrm{U})\) into \(\mathbf{C}\) . Then \(f\) is a distribution if and only if, for every compact subset \(\mathrm{K}\) of \(\mathrm{U}\) and for every family \((\mathrm{M}_{\alpha})_{\alpha \in \mathbf{N}^{n}}\) in \(\mathbf{R}_{+}\) , the linear form \(f\) is bounded on the set of functions \(\varphi \in \mathcal{C}_{\mathrm{K}}^{\infty}(\mathrm{U})\) such that, for every \(\alpha \in \mathbf{N}^{n}\) , we have

\[
\sup _ {x \in \mathrm{K}} | \partial^ {\alpha} \varphi (x) | \leqslant \mathrm{M} _ {\alpha}.
\]

Indeed, the space \(\mathcal{D}(\mathrm{U})\) is bornological and every bounded subset of \(\mathcal{D}(\mathrm{U})\) is contained in one of the bounded sets described in the statement.

Lemma 5. --- Let F be a filter on \(\mathcal{D}'(\mathrm{U})\) having a countable base or containing a simply bounded set. Then F converges to a distribution if and only if \(\langle f, \varphi \rangle\) converges according to F for every test function \(\varphi\) on U.

In particular, a sequence \((f_{m})_{m\in\mathbf{N}}\) of distributions converges to a distribution f if and only if, for every \(\varphi\in\mathcal{D}(\mathrm{U})\) the sequence \((\langle f_{m},\varphi\rangle)_{m\in\mathbf{N}}\) converges to \(\langle f,\varphi\rangle\) .

Since \(\mathcal{D}'(\mathrm{U})\) is a Montel space, hence barrelled; this follows from the Banach–Steinhaus theorem (EVT, III, p. 26, cor. 3 of th. 1).

Let V be an open set contained in U. The transpose of the canonical linear map from \(\mathcal{D}(\mathrm{V})\) into \(\mathcal{D}(\mathrm{U})\) is a continuous linear map from \(\mathcal{D}'(\mathrm{U})\) into \(\mathcal{D}'(\mathrm{V})\) , called the restriction of the distributions from U to V and denoted \(r_{\mathrm{VU}}\) , or sometimes \(r_{\mathrm{V,U}}\) .

One has \(r_{\mathrm{VV}} = 1_{\mathcal{D}'(\mathrm{V})}\) . If \(W \subset V \subset U\) are open subsets of U, then we have \(r_{\mathrm{WV}} \circ r_{\mathrm{VU}} = r_{\mathrm{WU}}\) . In other words, if \(\mathcal{T}\) denotes the topology on U, the projective system \(\underline{\mathcal{D}'}(\mathrm{U}) = ((\mathcal{D}'(\mathrm{V}))_{\mathrm{V} \in \mathcal{T}}, (r_{\mathrm{WV}}))\) is a presheaf on U with values in the species of structure of locally convex topological vector spaces (TA, I, p. 42, def. 1 and p. 66, §10).

PROPOSITION 5. --- The presheaf \(\underline{\mathcal{D}}^{\prime}(\mathrm{U})\) is a sheaf.

Recall (TA, I, p. 43, def. 2) that this means that for every open subset V of U and every open covering \((\mathrm{V}_{i})_{i\in\mathrm{I}}\) of V, the following conditions are satisfied:

\begin{enumerate}
\def\labelenumi{(\roman{enumi})}
\item
  The map \((r_{\mathrm{V}_i\mathrm{V}})_{i\in \mathrm{I}}\colon \mathcal{D}'(\mathrm{V})\to \prod_{i\in \mathrm{I}}\mathcal{D}'(\mathrm{V}_i)\) is injective;
\item
  For every family \((f_i)\in \prod_{i\in \mathrm{I}}\mathcal{D}'(\mathrm{V}_i)\) such that
\end{enumerate}

\[
r _ {(\mathrm{V} _ {i} \cap \mathrm{V} _ {i ^ {\prime}}) \mathrm{V} _ {i}} (f _ {i}) = r _ {(\mathrm{V} _ {i} \cap \mathrm{V} _ {i ^ {\prime}}) \mathrm{V} _ {i ^ {\prime}}} (f _ {i ^ {\prime}})
\]

for every pair \((i, i') \in \mathrm{I} \times \mathrm{I}\) , there exists a distribution \(f \in \mathcal{D}'(\mathrm{V})\) such that for every \(i \in I\) , one has \(r_{\mathrm{V}_{i}\mathrm{V}}(f) = f_{i}\) .

Let V be an open subset of U and \(\mathcal{V} = (\mathrm{V}_{i})_{i \in \mathrm{I}}\) an open covering of V. Let \((\mathrm{W}_{j})_{j \in \mathrm{J}}\) be an open covering of V that is locally finite, finer than \(\mathcal{V}\), and consisting of relatively compact parts (lemma 2 of IV, p. 201). Fix a partition of unity \((\varphi_{j})_{j \in \mathrm{J}}\) subordinate to the covering \((\mathrm{W}_{j})_{j \in \mathrm{J}}\) such that the support of \(\varphi_{j}\) is contained in \(W_{j}\) for all \(j \in J\) (lemma 1 of IV, p. 196).

Let \(\varphi\in\mathcal{D}(V)\) . Since the set of \(j\in J\) such that \(W_{j}\) meets the support of \(\varphi\) is finite (TG, I, §9, no. 1, p. 59), we have

\[
\varphi = \sum_ {j \in \mathrm{J}} \varphi \varphi_ {j}
\]

where the sum has only finitely many nonzero terms.

Let us prove (i). Suppose that \(f \in \mathcal{D}'(\mathrm{V})\) satisfies \(r_{\mathrm{V}_i\mathrm{V}}(f) = 0\) for all \(i \in \mathrm{I}\) . This means that \(\langle f, \varphi \rangle = 0\) for every test function \(\varphi\) whose support is contained in one of the open sets \(\mathrm{V}_i\) . This is a fortiori true if the support of \(\varphi\) is contained in one of the open sets \(\mathrm{W}_j\) . But then, for every \(\varphi \in \mathcal{D}(\mathrm{V})\) , we have

\[
\langle f, \varphi \rangle = \langle f, \sum_ {j \in \mathrm{J}} \varphi \varphi_ {j} \rangle = \sum_ {j \in \mathrm{J}} \langle f, \varphi \varphi_ {j} \rangle = 0,
\]

hence f = 0.

Let us prove (ii). Let \((f_{i})_{i\in\mathrm{I}}\) be a family such that \(f_{i}\in\mathcal{D}^{\prime}(\mathrm{V}_{i})\) for all \(i\in I\) and \(r_{\mathrm{V}_{i}\cap\mathrm{V}_{i^{\prime}},\mathrm{V}_{i}}(f_{i})=r_{\mathrm{V}_{i}\cap\mathrm{V}_{i^{\prime}},\mathrm{V}_{i^{\prime}}}(f_{i^{\prime}})\) for all i and \(i^{\prime}\) in I. Let \(\iota:J\to I\) be a map such that \(\mathrm{W}_{j}\subset\mathrm{V}_{\iota(j)}\) for all \(j\in J\) . For every \(j\in J\) set \(\widetilde{f}_{j}=r_{\mathrm{W}_{j},\mathrm{V}_{\iota(j)}}(f_{\iota(j)})\) . We then have

\[
\begin{array}{r l} & r _ {\mathrm{W} _ {j} \cap \mathrm{W} _ {j ^ {\prime}}, \mathrm{W} _ {j}} (\widetilde {f} _ {j}) = r _ {\mathrm{W} _ {j} \cap \mathrm{W} _ {j ^ {\prime}}, \mathrm{W} _ {j}} (r _ {\mathrm{W} _ {j}, \mathrm{V} _ {\iota (j)}} (f _ {\iota (j)})) \\ & \qquad = r _ {\mathrm{W} _ {j} \cap \mathrm{W} _ {j ^ {\prime}}, \mathrm{V} _ {\iota (j)}} (f _ {\iota (j)}) \\ & \qquad = r _ {\mathrm{W} _ {j} \cap \mathrm{W} _ {j ^ {\prime}}, \mathrm{V} _ {\iota (j)} \cap \mathrm{V} _ {\iota (j ^ {\prime})}} \circ r _ {\mathrm{V} _ {\iota (j)} \cap \mathrm{V} _ {\iota (j ^ {\prime})}, \mathrm{V} _ {\iota (j)}} (f _ {\iota (j)}). \end{array}
\]

By exchanging the roles of j and j' and noting that

\[
r _ {\mathrm{V} _ {\iota (j)} \cap \mathrm{V} _ {\iota (j ^ {\prime})}, \mathrm{V} _ {\iota (j)}} (f _ {\iota (j)}) = r _ {\mathrm{V} _ {\iota (j)} \cap \mathrm{V} _ {\iota (j ^ {\prime})}, \mathrm{V} _ {\iota (j ^ {\prime})}} (f _ {\iota (j ^ {\prime})}),
\]

by hypothesis, we deduce that

\[
r _ {\mathrm{W} _ {j} \cap \mathrm{W} _ {j ^ {\prime}}, \mathrm{W} _ {j}} (\widetilde {f} _ {j}) = r _ {\mathrm{W} _ {j} \cap \mathrm{W} _ {j ^ {\prime}}, \mathrm{W} _ {j}} (\widetilde {f} _ {j ^ {\prime}})\tag{1}
\]

for all \((j, j') \in \mathrm{J}^{2}\) .

For \(\varphi \in \mathcal{D}(\mathrm{V})\) , set

\[
\lambda (\varphi) = \sum_ {j \in \mathrm{J}} \langle \widetilde {f} _ {j}, \varphi \varphi_ {j} | \mathrm{W} _ {j} \rangle ,
\]

where the sum is finite since only finitely many terms can be nonzero.

The map \(\lambda\) is a linear form on \(\mathcal{D}(\mathrm{V})\) . Let us prove that it is a distribution. Let B be a bounded subset of \(\mathcal{D}(\mathrm{V})\) and let K be a compact subset of V such that \(B \subset \mathcal{C}_{K}^{\infty}(V)\) . By TG, I, §9, no. 1, p. 59, there exists a finite subset J' of J such that

\[
\lambda (\varphi) = \sum_ {j \in J ^ {\prime}} \left\langle \widetilde {f} _ {j}, \varphi \varphi_ {j} \mid W _ {j} \right\rangle
\]

for all \(\varphi\in\mathrm{B}\) . Since \(\widetilde{f}_{j}\) is a distribution for every \(j\in\mathrm{J}^{\prime}\) and that \(\mathcal{D}(\mathrm{V})\) is a topological algebra, it follows that \(\lambda\) is bounded on B, whence the result (Lemma 4).

Let \(j \in J\) and let \(\varphi \in \mathcal{D}(V)\) have support contained in \(W_{j}\) . We then have

\[
\langle \lambda , \varphi \rangle = \sum_ {j ^ {\prime} \in \mathrm{J}} \left\langle \widetilde {f} _ {j ^ {\prime}}, \varphi \varphi_ {j ^ {\prime}} \right| \mathrm{W} _ {j} \rangle = \sum_ {j ^ {\prime} \in \mathrm{J}} \left\langle \widetilde {f} _ {j}, \varphi \varphi_ {j ^ {\prime}} \right| \mathrm{W} _ {j} \rangle
\]

according to (1), since \(\varphi\varphi_{j'}\) has support contained in \(W_j \cap W_{j'}\) . Consequently

\[
\langle \lambda , \varphi \rangle = \langle \widetilde {f} _ {j}, \sum_ {j ^ {\prime} \in \mathrm{J}} \varphi \varphi_ {j ^ {\prime}} | \mathrm{W} _ {j} \rangle = \langle \widetilde {f} _ {j}, \varphi | \mathrm{W} _ {j} \rangle ,
\]

whence \(r_{\mathrm{W}_j\mathrm{V}}(\lambda) = \widetilde{f}_j\) for all \(j\in \mathbf{J}\) .

Let \(i \in I\) . Let us finally prove that the restriction of \(\lambda\) to \(V_i\) coincides with \(f_i\) . By condition (i), applied to the covering of \(V_i\) by the open sets \(W_j\) , it suffices to verify that for every \(j \in J\) the restriction of \(\lambda\) to \(V_i \cap W_j\) coincides with that of \(f_i\) . By the preceding, it remains to verify that the restriction of \(f_i\) to \(V_i \cap W_j\) is that of \(\widetilde{f}_j\) . Now, we have

\[
\begin{array}{c} r _ {\mathrm{V} _ {i} \cap \mathrm{W} _ {j}, \mathrm{V} _ {i}} (f _ {i}) = r _ {\mathrm{V} _ {i} \cap \mathrm{W} _ {j}, \mathrm{V} _ {i} \cap \mathrm{V} _ {\iota (j)}} (r _ {\mathrm{V} _ {i} \cap \mathrm{V} _ {\iota (j)}, \mathrm{V} _ {i}} (f _ {i})) \\ = r _ {\mathrm{V} _ {i} \cap \mathrm{W} _ {j}, \mathrm{V} _ {i} \cap \mathrm{V} _ {\iota (j)}} (r _ {\mathrm{V} _ {i} \cap \mathrm{V} _ {\iota (j)}, \mathrm{V} _ {\iota (j)}} (f _ {\iota (j)})) = r _ {\mathrm{V} _ {i} \cap \mathrm{W} _ {j}, \mathrm{V} _ {\iota (j)}} (f _ {\iota (j)}), \end{array}
\]

where we used the hypothesis concerning the family \((f_{i})\) . But then

\[
r _ {\mathrm{V} _ {i} \cap \mathrm{W} _ {j}, \mathrm{V} _ {\iota (j)}} (f _ {\iota (j)}) = r _ {\mathrm{V} _ {i} \cap \mathrm{W} _ {j}, \mathrm{W} _ {j}} (r _ {\mathrm{W} _ {j}, \mathrm{V} _ {\iota (j)}} (f _ {\iota (j)})) = r _ {\mathrm{V} _ {i} \cap \mathrm{W} _ {j}, \mathrm{W} _ {j}} (\widetilde {f} _ {j}),
\]

which allows us to conclude.

